\section{Static Storage and Conversion Sequencing}
The controller uses the CMOS cells in Fig.~\ref{fig:logic-cells}, not ideal
logic sources. Figure~\ref{fig:logic-latch} expands its resettable storage
hierarchy. With $en=1$, the input transmission gate conducts and the feedback
gate opens. With $en=0$, the feedback gate retains the state. Pulling
$nrst=\overline{RST}$ low forces the NAND output $qb$ high and the final
output $q$ low. The formal $rst$ latch port is retained in the hierarchy;
the reset-to-zero action is applied through $nrst$.

\begin{figure*}[!t]
\centering
\resizebox{.97\textwidth}{!}{% Direct terminal order is D G S B; dimensions are nominal W L in metres.
% @mos sar_inv Mp PMOS y a vdd vdd 1e-6 30e-9
% @mos sar_inv Mn NMOS y a vss vss .5e-6 30e-9
% @mos sar_nand2 Mp0 PMOS y a vdd vdd 1e-6 30e-9
% @mos sar_nand2 Mp1 PMOS y b vdd vdd 1e-6 30e-9
% @mos sar_nand2 Mn0 NMOS y a ns vss 1e-6 30e-9
% @mos sar_nand2 Mn1 NMOS ns b vss vss 1e-6 30e-9
% @mos sar_tg Mn NMOS y en a vss .5e-6 30e-9
% @mos sar_tg Mp PMOS y enb a vdd 1e-6 30e-9
\begin{circuitikz}
% (a) CMOS inverter; named terminals follow the exported device order.
\node[pmos] (ip) at (1.65,4.0) {};
\node[nmos] (in) at (1.65,1.65) {};
\draw (ip.S) -- (1.65,5.05) node[above] {$V_{\mathrm{DD}}$};
\draw (in.S) -- (1.65,.70) node[ground] {};
\draw (ip.D) -- (in.D);
\draw (ip.G) -- (.35,4.0) -- (.35,1.65) -- (in.G);
\draw (.35,2.8) node[adc dot] {} -- (-.2,2.8) node[left] {$a$};
\draw (1.65,2.8) node[adc dot] {} -- (2.7,2.8) node[right] {$y$};
\node[right] at (ip.D) {$M_p$};
\node[right] at (in.D) {$M_n$};
\node[align=center] at (1.4,-1.0) {(a) Inverter\\$W_p/W_n=1/0.5\;\mu\mathrm{m}$};
% (b) NAND2: parallel PMOS pull-up and series NMOS pull-down.
\node[pmos] (np0) at (5.4,4.0) {};
\node[pmos,xscale=-1] (np1) at (7.8,4.0) {};
\node[nmos] (nn0) at (6.6,2.15) {};
\node[nmos] (nn1) at (6.6,.65) {};
\draw (np0.S) -- (5.4,5.05) -- (7.8,5.05) -- (np1.S);
\node[above] at (6.6,5.05) {$V_{\mathrm{DD}}$};
\draw (np0.D) -- (5.4,3.1) -- (7.8,3.1) -- (np1.D);
\draw (6.6,3.1) node[adc dot] {} -- (nn0.D);
\draw (7.8,3.1) node[adc dot] {} -- (8.55,3.1) node[right] {$y$};
\draw (nn0.S) -- (nn1.D);
\draw (6.6,1.38) node[adc dot] {} -- (7.15,1.38) node[right] {$ns$};
\draw (nn1.S) -- (6.6,-.15) node[ground] {};
\draw (np0.G) -- ++(-.35,0) node[left] {$a$};
\draw (np1.G) -- ++(.35,0) node[right] {$b$};
\draw (nn0.G) -- ++(-.35,0) node[left] {$a$};
\draw (nn1.G) -- ++(-.35,0) node[left] {$b$};
\node[right] at (np0.D) {$M_{p0}$};
\node[left] at (np1.D) {$M_{p1}$};
\node[right] at (7.0,2.15) {$M_{n0}$};
\node[right] at (7.0,.65) {$M_{n1}$};
\node[align=center] at (6.6,-1.0) {(b) NAND2\\All four widths: $1\;\mu\mathrm{m}$};
% (c) Parallel complementary pass devices; a is the declared source.
\node[nmos,rotate=-90] (tn) at (12.35,3.65) {};
\node[pmos,rotate=90] (tp) at (12.35,1.95) {};
\draw (tn.S) -- (11.2,3.65) -- (11.2,1.95) -- (tp.S);
\draw (tn.D) -- (13.5,3.65) -- (13.5,1.95) -- (tp.D);
\draw (11.2,2.8) node[adc dot] {} -- (10.6,2.8) node[left] {$a$};
\draw (13.5,2.8) node[adc dot] {} -- (14.1,2.8) node[right] {$y$};
\draw (tn.G) -- ++(0,.4) node[above] {$en$};
\draw (tp.G) -- ++(0,-.4) node[below] {$enb$};
\node[above left] at (tn.S) {$M_n$};
\node[below left] at (tp.S) {$M_p$};
\node[align=center] at (12.35,-1.0) {(c) Transmission gate\\$W_n/W_p=0.5/1\;\mu\mathrm{m}$};
\node[align=center] at (7.0,-2.0) {All $L=30\;\mathrm{nm}$. NMOS bulk: $V_{\mathrm{SS}}$; PMOS bulk: $V_{\mathrm{DD}}$.\\Equal labels denote the same net; dots mark junctions. Ground denotes $V_{\mathrm{SS}}$.};
\end{circuitikz}
}
\caption{Source-derived transistor implementations of the standard inverter,
NAND2 and transmission gate. Each cell uses its own local net names. NMOS
bulks connect to $V_{\mathrm{SS}}$ and PMOS bulks to $V_{\mathrm{DD}}$;
the three-terminal symbols omit these common body wires. Dimensions are
nominal requested geometry, not extracted dimensions or a replacement model card.}
\label{fig:logic-cells}
\end{figure*}

\begin{table}[!t]\centering
\caption{Logic-device terminal order $(D,G,S,B)$}
\label{tab:logic-terminals}
\begin{tabular}{lll}\toprule
Cell & Device & $(D,G,S,B)$\\\midrule
Inverter & $M_p$ & $(y,a,vdd,vdd)$\\
 & $M_n$ & $(y,a,vss,vss)$\\
NAND2 & $M_{p0}$ & $(y,a,vdd,vdd)$\\
 & $M_{p1}$ & $(y,b,vdd,vdd)$\\
 & $M_{n0}$ & $(y,a,ns,vss)$\\
 & $M_{n1}$ & $(ns,b,vss,vss)$\\
TG & $M_n$ & $(y,en,a,vss)$\\
 & $M_p$ & $(y,enb,a,vdd)$\\\bottomrule
\end{tabular}
\end{table}

\begin{figure*}[!t]
\centering
\resizebox{.97\textwidth}{!}{% @instance sar_latchr0 Xin sar_tg vdd vss d x en enb
% @instance sar_latchr0 Xfb sar_tg vdd vss q x enb en
% @instance sar_latchr0 Xg sar_nand2 vdd vss x nrst qb
% @instance sar_latchr0 Xo sar_inv vdd vss qb q
% @instance sar_pdffr0 Xmaster sar_latchr0 vdd vss d phm phmb rst nrst mq
% @instance sar_pdffr0 Xslave sar_latchr0 vdd vss mq phs phsb rst nrst q
\begin{circuitikz}
% Actual latch hierarchy, whose TG and gate implementations are shown separately.
\node[adc block,minimum width=1.8cm,minimum height=.95cm] (xin) at (2.4,3.0) {TG\\$X_{\mathrm{in}}$};
\node[nand port] (ng) at (6.7,3.0) {};
\node[not port] (iv) at (9.0,3.0) {};
\node[adc block,minimum width=1.8cm,minimum height=.95cm] (xfb) at (6.7,.75) {TG\\$X_{\mathrm{fb}}$};
\draw (.1,3.0) node[left] {$d$} -- (xin.west);
\draw (xin.east) -- (4.2,3.0) |- (ng.in 1);
\node[above] at (4.2,3.0) {$x$};
\draw (5.0,2.25) node[left] {$nrst$} -| (ng.in 2);
\draw (ng.out) -- (iv.in);
\node[above] at (7.8,3.0) {$qb$};
\draw (iv.out) -- (11.1,3.0) node[right] {$q$};
\draw (10.25,3.0) node[adc dot] {} -- (10.25,.75) -- (xfb.east);
\draw (xfb.west) -- (4.2,.75) -- (4.2,3.0) node[adc dot] {};
\draw (xin.north) -- ++(0,.42) node[above] {$en$};
\draw (xin.south) -- ++(0,-.42) node[below] {$enb$};
\draw (xfb.north) -- ++(0,.42) node[above] {$enb$};
\draw (xfb.south) -- ++(0,-.42) node[below] {$en$};
\node[above] at (ng.north) {$X_g$};
\node[above] at (iv.north) {$X_o$};
\node[align=left,anchor=west] at (11.8,1.6) {$q=x\,nrst$\\$nrst=\overline{RST}$\\10 MOS per latch};
% DFF is a hierarchy of two physical latches, not an ideal behavioral device.
\node[adc block,minimum width=2.35cm,minimum height=1.0cm] (ml) at (4.0,-2.0) {Master latch\\$X_{\mathrm{master}}$};
\node[adc block,minimum width=2.35cm,minimum height=1.0cm] (sl) at (8.4,-2.0) {Slave latch\\$X_{\mathrm{slave}}$};
\draw (1.7,-2.0) node[left] {$d$} -- (ml.west);
\draw (ml.east) -- (sl.west) node[midway,above] {$mq$};
\draw (sl.east) -- (10.7,-2.0) node[right] {$q$};
\draw (ml.north) -- ++(0,.35) node[above] {$phm,\ phmb$};
\draw (sl.north) -- ++(0,.35) node[above] {$phs,\ phsb$};
\draw (ml.south) -- ++(0,-.35) node[below] {$RST,\ nrst$};
\draw (sl.south) -- ++(0,-.35) node[below] {$RST,\ nrst$};
\node[align=left,anchor=west] at (11.8,-2.0) {\texttt{sar\_pdffr0}\\20 MOS per DFF};
\node[align=center] at (7.2,-3.75) {Each TG box denotes a complementary two-MOS pass gate.\\Supplies $V_{\mathrm{DD}},V_{\mathrm{SS}}$ are common to every cell.};
\end{circuitikz}
}
\caption{Actual reset-to-zero latch hierarchy and its two-latch DFF composition.
The input TG uses $(en,enb)$ and the feedback TG uses $(enb,en)$.
Each latch contains ten MOS instances; the DFF contains twenty, excluding
the shared phase generator and reset inverter.}
\label{fig:logic-latch}
\end{figure*}

The five Johnson stages obey
\begin{equation}
(j_0,j_1,j_2,j_3,j_4)^+
  =(\overline{j_4},j_0,j_1,j_2,j_3).
\end{equation}
Reset initializes all five state bits to zero. There are ten decoded slots,
not ten physical state flip-flops. The first two slots acquire the input;
the following eight select decisions from bit~7 down to bit~0
(Table~\ref{tab:johnson-slots}). Eight bit latches and eight output latches
complete the storage shown in Fig.~\ref{fig:controller}.
The physical phase generator uses two cross-inhibited NOR paths with
buffers: $mraw=\overline{clkb+phs}$ and
$sraw=\overline{clk+phm}$. Its finite gate delays create the actual
master/slave phase separation; the drawing does not replace it with ideal
nonoverlap clocks.

\begin{table}[!t]\centering
\caption{Johnson states and conversion operations}
\label{tab:johnson-slots}
\begin{tabular}{cll}\toprule
Slot & $j_0j_1j_2j_3j_4$ & Operation / select $s_k$\\\midrule
0 & 00000 & Acquire; output latch enabled by $phm$\\
1 & 10000 & Acquire\\
2 & 11000 & Bit 7: $j_1\overline{j_2}$\\
3 & 11100 & Bit 6: $j_2\overline{j_3}$\\
4 & 11110 & Bit 5: $j_3\overline{j_4}$\\
5 & 11111 & Bit 4: $j_0j_4$\\
6 & 01111 & Bit 3: $\overline{j_0}j_1$\\
7 & 00111 & Bit 2: $\overline{j_1}j_2$\\
8 & 00011 & Bit 1: $\overline{j_2}j_3$\\
9 & 00001 & Bit 0: $\overline{j_3}j_4$\\\bottomrule
\end{tabular}
\end{table}

\begin{figure*}[!t]
\centering
\resizebox{.97\textwidth}{!}{% Source-derived hierarchy; all five storage stages are sar_pdffr0 cells.
% @instance sar_control Xjohnson0 sar_pdffr0 vdd vss j4b phm phmb phs phsb RST nrst j0
% @instance sar_control Xjohnson1 sar_pdffr0 vdd vss j0 phm phmb phs phsb RST nrst j1
% @instance sar_control Xjohnson2 sar_pdffr0 vdd vss j1 phm phmb phs phsb RST nrst j2
% @instance sar_control Xjohnson3 sar_pdffr0 vdd vss j2 phm phmb phs phsb RST nrst j3
% @instance sar_control Xjohnson4 sar_pdffr0 vdd vss j3 phm phmb phs phsb RST nrst j4
% @instance sar_control Xjohnson_inv4 sar_inv vdd vss j4 j4b
% @instance sar_control Xq0_decode sar_and2 vdd vss j0b j4b q0
% @instance sar_control Xtrack sar_nor2 vdd vss j1 j4 track
% @instance sar_control Xcompare_clock sar_and2 vdd vss phm ntrack cmpclk
% @instance sar_control Xoutput_enable_n sar_nand2 vdd vss q0 phm oenb
% @instance sar_control Xoutput_enable sar_inv_hi vdd vss oenb oen
\begin{circuitikz}
\foreach \i/\x in {0/1.4,1/4.0,2/6.6,3/9.2,4/11.8}{
  \node[adc block,minimum width=1.65cm,minimum height=1.1cm] (ff\i) at (\x,3.3) {DFF\\$j_{\i}$};
  \draw (ff\i.south) -- ++(0,-.35);
}
\draw (ff0.east) -- (ff1.west) node[midway,above] {$j_0$};
\draw (ff1.east) -- (ff2.west) node[midway,above] {$j_1$};
\draw (ff2.east) -- (ff3.west) node[midway,above] {$j_2$};
\draw (ff3.east) -- (ff4.west) node[midway,above] {$j_3$};
\node[not port] (fb) at (14.0,3.3) {};
\draw (ff4.east) -- (fb.in) node[midway,above] {$j_4$};
\draw (fb.out) -- (15,3.3) -- (15,5.0) -- (-.1,5.0) -- (-.1,3.3) -- (ff0.west);
\node[above] at (7.4,5.0) {$j_{4b}=\overline{j_4}$};
\draw (.4,2.4) -- (12.8,2.4);
\foreach \x in {1.4,4.0,6.6,9.2,11.8}{\draw (\x,2.75) -- (\x,2.4) node[adc dot] {};}
\node[below] at (6.6,2.4) {Common phase/reset ports: $phm,phmb,phs,phsb,RST,nrst$};
\node[adc block,minimum width=4.7cm,minimum height=2.0cm,align=left] at (2.35,.25)
  {State decoding\\$q_0=\overline{j_0}\,\overline{j_4}$\\$track=\overline{j_1+j_4}$\\$cmpclk=phm\,\overline{track}$};
\node[adc block,minimum width=4.0cm,minimum height=2.0cm,align=left] at (7.2,.25)
  {Bit storage, $k=7\ldots0$\\$ben_k=s_k\,phm$\\$b_k\leftarrow cmpQ$\\8 physical reset latches};
\node[adc block,minimum width=4.65cm,minimum height=2.0cm,align=left] at (12.3,.25)
  {Output storage\\$oen=q_0\,phm$\\$D_k\leftarrow b_k$\\8 physical reset latches};
\node[align=center] at (7.25,-1.55) {Five master/slave DFFs form ten decoded slots. Blocks show cell hierarchy, not ideal simulation elements.};
\end{circuitikz}
}
\caption{Five-stage Johnson controller and the actual state-decoding relations.
Every DFF expands to the physical master/slave hierarchy in
Fig.~\ref{fig:logic-latch}. The horizontal phase/reset line is a labeled bundle
of six distinct nets, not a short between those signals. The decoded bit and
output storage blocks denote repeated physical latches. The active control
hierarchy contains 424 MOS instances.}
\label{fig:controller}
\end{figure*}
